% !PRINTED 417
\weberpage{417}
It is also shown there how the general equation of degree $5$ can be transformed into one corresponding to these assumptions.

Under the assumption (4), the subsistence of the relations (1) has no influence on the conclusions of \S.\,108, and we then have the following theorems.

Let $x_0, x_1, \ldots, x_{n-1}$ be a system of variables which is subject to the relations $A=0$, $B=0$, \ldots\ but otherwise independent, and let $\mathfrak{A}$ be a permutation group of the $x$ whose permutations leave the functions $A, B, \ldots$ unchanged.

\gesperrt{If the equation whose roots are the quantities $x$ is to possess a resolvent}
\weq{6}{F(u,z) = 0}
\gesperrt{which, apart from numerical coefficients, contains only one parameter that is a function of the $x$ unchanged by the group $\mathfrak{A}$, then there must exist a function $\vartheta$ of the variables $x_0, \ldots, x_{n-1}$ which, under the permutations $\alpha$ of $\mathfrak{A}$ (having regard to the relations $A=0$, $B=0$, \ldots), undergoes linear substitutions}
\weq{7}{\vartheta_\alpha = \dfrac{a\vartheta+b}{c\vartheta+d},}
\gesperrt{in which $a, b, c, d$ are numerical coefficients. The roots of the equation (6) are rationally expressible in terms of $\vartheta$, and consequently $\vartheta$ is also the root of an equation with one parameter}
\weq{8}{\Phi(\vartheta, z) = 0.}

If (6) is a total resolvent, which always occurs when the group $\mathfrak{A}$ is simple, then the $x$ too can be expressed rationally in terms of $\vartheta$ and the functions belonging to the group $\mathfrak{A}$.

Up to this point, then, there is complete agreement with the results of \S.\,108.

But the restriction upon which we came in \S.\,109 will here no longer necessarily arise, because now it is no longer possible, as in the case of completely independent variables, to conclude from the equality of two fractional functions the agreement of numerator and denominator.

The linear substitutions (7) form a group $P$, (simply or multiply) isomorphic with $\mathfrak{A}$, which with one of our
%CONT>>
